\documentclass[11pt]{article}
\usepackage[utf8]{inputenc}
\usepackage{amsmath,amssymb}
\usepackage{xcolor}
\usepackage{array}
\usepackage{natbib}
\usepackage{graphicx}
\DeclareGraphicsExtensions{.pdf,.png,.jpg}
\usepackage{float}
\usepackage[top=3cm, bottom=3cm, right=3cm, left=3cm]{geometry}
\usepackage{url}
\usepackage{booktabs}
\usepackage{enumitem}
\usepackage{authblk}
\usepackage{multirow}

\newcommand{\placeholder}[1]{\textcolor{gray}{\textit{[#1]}}}
\newcommand{\todo}[1]{\textcolor{red}{\textbf{[TODO: #1]}}}

\title{Do LLMs Interact Like Humans? \\ Endogenous Interaction Dynamics Under Roleplay \\ and Prediction Conditions Across Nine Social Domains}

\author[1]{C. Ben Gibson}
\affil[1]{\placeholder{Department, Institution, City, Country}}

\begin{document}

\maketitle

\noindent\textbf{Corresponding author:} C. Ben Gibson,
\placeholder{email@institution.edu}

\begin{abstract}
Large language models (LLMs) are increasingly deployed as social simulation
agents, yet it remains unclear whether they reproduce the structural dynamics
that govern human interaction sequences. We evaluate two LLMs---Claude Haiku
4.5 and Llama 3.1 70B---under two conditions that differ in how much
structural scaffolding is provided: a \emph{roleplay} condition, in which each
LLM is assigned a sender identity and selects only a receiver (with sender
activation governed by an empirically fitted relational event model), and a
\emph{predict} condition, in which the LLM observes the empirical event
history and predicts both the next sender and receiver with no structural
scaffolding. We compare the resulting sequences against human empirical
baselines on nine endogenous statistics from the relational event model
literature, across 27 networks spanning nine social domains. The two
conditions reveal sharply different failure modes. In the roleplay condition,
both LLMs systematically \emph{over-produce} history-dependent statistics:
too much dyadic repetition (Claude $+97\%$, Llama $+44\%$), too much
reverse-dyad recency concentration ($+61\%$ for both), and dramatically
too much contact breadth ($+189\%$ and $+227\%$), while under-producing
interaction-order violations ($-26\%$ for both). In the predict condition,
Claude \emph{reverses} these biases, sharply under-producing the same
history-dependent statistics (SndInitiation $-55\%$, Repetition $-80\%$,
RecRecencyRank $-72\%$) while over-producing norm violations ($+44\%$),
yielding lower overall RMSE (0.33) than in roleplay (1.66) but through
cancellation rather than accuracy. Llama in the predict condition is
strikingly well-calibrated (RMSE 0.19), reproducing most statistics close
to human baselines without identity assignment. The contrast between
conditions---over-endogeneity under roleplay, under-endogeneity under
prediction for Claude, and near-accuracy for Llama---has implications for
how LLMs should be deployed in social simulation and for what these models
have internalized about human interaction dynamics.
\end{abstract}

\noindent\textbf{Keywords:} large language models; social simulation; relational
event models; interaction dynamics; agent-based modeling; endogenous statistics;
roleplay; prediction

\newpage

%------------------------------------------------------------------
\section{Introduction}

A growing body of work uses large language models (LLMs) as computational
proxies for human social behavior. LLM agents have been deployed to simulate
opinion dynamics \citep{argyle2023out}, political deliberation
\citep{manning2024automated}, jury decision-making, and economic games
\citep{horton2023large,aher2023using}, with researchers variously concluding that
LLMs are plausible synthetic survey respondents, strategic actors, and social
reasoners. A parallel literature in organizational science and computational
social science employs LLMs as interacting agents in network simulations,
leveraging their world knowledge to generate scenario-conditioned behavior at
scale. Both applications rest on an implicit empirical claim: that LLM behavior
in the relevant domain resembles human behavior closely enough to serve as a
valid substitute.

This claim has received less direct empirical scrutiny than the enthusiasm
surrounding LLM social simulation might suggest. Most evaluations focus on
macro-level outcomes---whether aggregate opinion distributions match survey
data, whether strategic equilibria are reproduced, whether argumentation
patterns resemble human debates. Micro-level sequencing---the fine-grained,
moment-to-moment dynamics of who interacts with whom, in what order, and
how that pattern evolves over time---has received much less attention. Yet
these dynamics are precisely what social network researchers study, and they
are what matters most for simulation applications where the quality of
simulated trajectories, not just their endpoints, is the object of interest.

We address this gap by evaluating LLM-generated interaction sequences against
human empirical baselines on nine endogenous statistics drawn from the
relational event model (REM) literature \citep{butts_sm_2008,gibson2025four},
across 27 networks spanning nine social domains. Crucially, we evaluate LLMs
under two conditions that differ in the degree of structural scaffolding
provided.

In the \emph{roleplay} condition, each LLM is assigned a specific sender
identity and asked to select a receiver. Sender activation at each step is
governed by the hazard surface of an empirically fitted REM, concentrating
event production on actors whose prior interaction history places them at
elevated hazard. The LLM's task is narrow: given your identity and recent
history, who do you target next?

In the \emph{predict} condition, the LLM receives the same empirical event
history but is asked to predict both the next sender and the next receiver,
with no REM structural scaffolding. This is a harder and qualitatively
different task: the LLM must represent the full distribution over
(sender, receiver) pairs and draw from an implicit model of
who acts and who they target.

The two conditions operationalize a conceptual distinction between
\emph{identity-conditioned behavior}---acting as a specific social agent---and
\emph{social forecasting}---predicting what will happen in a social sequence.
We hypothesize that identity assignment will amplify endogenous
history-responsiveness, because the roleplay frame invites the LLM to adopt
the perspective of a persistent actor with established relationships and
obligations. The predict frame, by contrast, requires the LLM to represent
population-level tendencies rather than individual behavioral dispositions,
which may yield different calibration properties.

Our results confirm this hypothesis but with a more nuanced pattern than
anticipated. In the roleplay condition, both LLMs \emph{over-produce}
endogenous structure across most statistics, behaving as if the pull of
interaction history is stronger than it is for humans. In the predict
condition, Claude reverses this pattern and sharply \emph{under-produces}
history-dependent statistics, while Llama turns out to be nearly
well-calibrated overall. The contrast---over-endogeneity under roleplay,
under-endogeneity under prediction (for Claude), and implicit accuracy under
prediction (for Llama)---suggests that the two models have different
internal representations of social interaction dynamics, and that these
representations are expressed differently depending on whether the LLM is
asked to act or to forecast.

The remainder of the paper proceeds as follows. Section~\ref{sec:background}
reviews endogenous REM statistics and prior work on LLMs as social simulation
agents. Section~\ref{sec:methods} describes the data, conditions, and
evaluation metrics. Section~\ref{sec:results} reports results for both
conditions and their comparison. Section~\ref{sec:discussion} discusses
interpretations and implications.

%------------------------------------------------------------------
\section{Background}
\label{sec:background}

\subsection{Endogenous Structure in Relational Event Sequences}

The relational event model \citep{butts_sm_2008} represents directed social
interaction as a continuous-time point process in which the hazard that actor
$i$ directs an event to actor $j$ at time $t$ depends on the full prior event
history:
\[
  \lambda_{ijt} = \lambda_0(t)\exp\!\left(\boldsymbol{\beta}^\top
  \mathbf{s}_{ijt}\right),
\]
where $\mathbf{s}_{ijt}$ is a vector of endogenous statistics computed from
events up to time $t$. The framework has been applied to email and phone
communication \citep{vu2017relational}, hospital patient transfers
\citep{liang2014organizational}, organizational deliberation
\citep{kitts2017investigating}, gang conflict \citep{duxbury2021shining,%
gravel2023rivalries}, and international conflict \citep{renshaw2023modeling}.

\citet{gibson2025four} applied greedy forward model selection to 217 empirical
interaction sequences across nine social domains and identified four statistics
that appear in the first four selection steps in over 60\% of networks:
\textit{SndInitiation} (whether the sender has ever previously contacted the
receiver), \textit{IOviolation} (indicator of participation-shift violations),
\textit{Repetition} (sender's normalized prior send-share to the receiver), and
\textit{RecRecencyRank} (rank-based recency of the most recent reverse-dyad
event). These four terms, representing dyadic persistence, sequential ordering,
and reverse-dyad recency, constitute a robust endogenous baseline that
generalizes across ecological contexts from private messaging to gang violence.

An ecological distinction that shapes which endogenous statistics dominate is
the \emph{visibility} of interaction histories: whether participants can observe
each other's behavior \citep{gibson2025four}. In high-visibility settings
(classrooms, public conflict, Twitter), imitation and reverse-dyad recency are
more predictive because others' targeting behavior is observable. In
low-visibility settings (private email, phone, SMS), actors access only their
own interaction history, so dyadic persistence and sequential norms dominate.

\subsection{LLMs as Social Simulation Agents}

The use of LLMs as synthetic social agents has grown rapidly since the
release of instruction-tuned models. Early work established that LLMs can
reproduce statistical patterns in survey responses when prompted to adopt
specific demographic profiles \citep{argyle2023out}, though subsequent work
documented systematic biases and over-homogenization
\citep{santurkar2023whose,park2024generative}. In strategic interaction
settings, LLMs demonstrate economically rational behavior with characteristic
deviations from Nash equilibria \citep{horton2023large}. Generative agent
architectures \citep{park2023generative} have extended LLM simulation to
multi-agent settings where agents maintain memories and coordinate behavior
over extended sequences.

A growing body of work documents that persona and roleplay prompting
systematically distorts LLM behavior relative to neutral or direct prompting.
\citet{shanahan2023role} argue that LLMs in deep roleplay are best understood
as \emph{simulating} a character rather than becoming one, and identify
\emph{character capture}---the tendency for strong roleplay framing to
override the model's baseline values and behavioral dispositions---as a
central failure mode of identity-based prompting. Empirically,
\citet{cheng2023compost} find that LLMs roleplaying demographic groups
systematically exaggerate stereotypical traits, producing what they term
\emph{caricature}: the assigned persona amplifies group-typical tendencies
beyond what real members of that group would exhibit.
\citet{salewski2023context} similarly show that impersonating specific
personas shifts task performance and expressed opinions in
stereotype-consistent directions. Taken together, these results suggest that
identity assignment does not produce an accurate simulation of the target
actor; it produces a \emph{heightened} version of whatever tendencies the
model associates with that identity. \citet{tseng2024two} survey the
distinction between role-playing and personalization in LLMs, arguing that
the two paradigms have distinct failure modes that are frequently conflated
in evaluation.

An important distinction in LLM simulation is between \emph{identity-based}
prompting---assigning the model a specific actor's role and history---and
\emph{predictive} prompting---asking the model to forecast what a social
sequence will produce next. Prior work on LLM simulation of conflict
sequences found that sender roleplay substantially outperforms prediction on
REM statistics, with identity assignment (not merely history observation)
driving the accuracy improvement \citep{gibson2024llm}. That work evaluated
a single domain; here we extend the comparison across nine domains and
examine how the two framing conditions produce different patterns of
systematic bias.

%------------------------------------------------------------------
\section{Methods}
\label{sec:methods}

\subsection{Empirical Corpora}

The evaluation draws on 27 networks across nine social domains:

\begin{itemize}[leftmargin=2em]
  \item \textbf{War} (international conflict, MID dataset):
        directed conflict acts among states.
        High-visibility, conflictual. ($n=1$ network)
  \item \textbf{Classroom}: directed spoken interaction in a high-school
        classroom. High-visibility, affiliative. ($n=1$)
  \item \textbf{Gangs}: directed violent acts among urban street gangs.
        High-visibility, conflictual. ($n=5$)
  \item \textbf{Twitter}: directed replies and mentions.
        High-visibility, affiliative. ($n=1$)
  \item \textbf{Phone}: phone calls in a military command-and-control exercise.
        High-visibility, affiliative. ($n=1$)
  \item \textbf{Email}: university email sequences.
        Low-visibility, affiliative. ($n=5$)
  \item \textbf{ftf (face-to-face)}: directed gaze sequences during a group task.
        Low-visibility, affiliative. ($n=5$)
  \item \textbf{PMs (private messages)}: UCI online private message sequences.
        Low-visibility, affiliative. ($n=5$)
  \item \textbf{SMS}: short message service sequences.
        Low-visibility, affiliative. ($n=3$)
\end{itemize}

For each network, the empirical sequence serves as the human baseline
against which LLM-generated sequences are compared on nine endogenous
statistics.

\subsection{LLM Providers}

Two LLM providers are evaluated throughout:
\begin{itemize}[leftmargin=2em]
  \item \textbf{Claude Haiku 4.5} (Anthropic), via AWS Bedrock
  \item \textbf{Llama 3.1 70B} (Meta), via AWS Bedrock
\end{itemize}

\subsection{Experimental Conditions}

\paragraph{Roleplay condition.}
Each LLM is assigned the identity of a specific sender at each event step and
asked to select a receiver from the available actors, given the last 30 events
as context. Sender activation at each step is governed by the hazard surface
of an empirically fitted REM: at each tick, sender activation probabilities
$p_i(t)$ are derived by marginalizing the Cox partial hazard over receivers
and normalizing, concentrating activations on actors whose prior interaction
history places them at elevated hazard. The LLM's task is narrow: ``You are
[Actor $i$]. Given this history, who do you target next?'' This condition
provides structural scaffolding for the \emph{sender} dimension while leaving
the \emph{receiver} selection to the LLM.

\paragraph{Predict condition.}
The LLM receives the same empirical event history (last 30 events) at each
step but is asked to predict both the next sender and the next receiver,
with no REM structural scaffolding. The LLM sees the true historical sequence
up to each event step (\texttt{use\_empirical\_history = TRUE}) and is prompted
to name the next (sender, receiver) pair. This is a harder and qualitatively
different task: rather than acting as a specific agent, the LLM must forecast
the next event from a population-level model of interaction.

The contrast between conditions is summarized in
Table~\ref{tab:conditions}.

\begin{table}[H]
\centering
\caption{Summary of experimental conditions.}
\label{tab:conditions}
\begin{tabular}{lll}
\toprule
 & \textbf{Roleplay} & \textbf{Predict} \\
\midrule
LLM task           & Choose receiver only          & Choose sender and receiver \\
Sender determined by & REM hazard surface           & LLM \\
Identity assigned  & Yes (specific actor)          & No \\
History observed   & Last 30 events (LLM-generated) & Last 30 empirical events \\
REM scaffolding    & Yes (sender activation)       & None \\
Difficulty         & Lower                         & Higher \\
\bottomrule
\end{tabular}
\end{table}

\subsection{Endogenous Statistics}

Nine endogenous statistics from the REM candidate set of \citet{gibson2025four}
are evaluated (Table~\ref{tab:terms}). Each is normalized to $[0,1]$ within
sender at each event time.

\begin{table}[H]
\centering
\caption{Endogenous statistics evaluated, organized by conceptual group.}
\label{tab:terms}
\begin{tabular}{llp{7.2cm}}
\toprule
\textbf{Term} & \textbf{Group} & \textbf{Description} \\
\midrule
SndInitiation   & Forward-dyad history   & Binary: sender $i$ has ever previously contacted receiver $j$ \\
Repetition      & Forward-dyad history   & $i$'s normalized prior send-share to $j$ \\
RecRecencyRank  & Reverse-dyad history   & Rank-recency of the most recent $j \to i$ event \\
WeightedRec     & Reverse-dyad history   & Squartini weighted reciprocity ratio \\
Imitation       & Popularity/closure     & $j$'s normalized indegree from senders other than $i$ \\
OutDegSnd       & Popularity/closure     & $i$'s normalized contact breadth (sender outdegree) \\
NOTPSnd         & Popularity/closure     & Normalized transitive two-path count ($i \to k \to j$) \\
NITPSnd         & Popularity/closure     & Normalized cyclic two-path count ($j \to k \to i$) \\
IOviolation     & Sequential ordering    & Indicator of a participation-shift violation \\
\bottomrule
\end{tabular}
\end{table}

\subsection{Evaluation Metrics}

For each network and each statistic, bias is the difference between the
LLM-generated mean and the empirical mean:
\[
  \text{bias}_{k} = \bar{s}_{k}^{\text{LLM}} - \bar{s}_{k}^{\text{emp}}.
\]
Positive bias indicates the LLM over-produces the pattern; negative
indicates under-production. Root mean squared error (RMSE) across the nine
statistics is computed per network:
\[
  \text{RMSE} = \sqrt{\frac{1}{9}\sum_{k=1}^{9}
    \!\left(\bar{s}_{k}^{\text{LLM}} - \bar{s}_{k}^{\text{emp}}\right)^{\!2}}.
\]

%------------------------------------------------------------------
\section{Results}
\label{sec:results}

\subsection{Overall Calibration}

Table~\ref{tab:overall} reports mean bias, mean absolute bias, and
RMSE for each condition and model. Several patterns stand out immediately.

\begin{table}[H]
\centering
\caption{Overall calibration. Mean bias (positive = over-production relative to
humans), mean absolute bias, and RMSE. Averages are across all 9 statistics
and 27 networks.}
\label{tab:overall}
\begin{tabular}{llccc}
\toprule
\textbf{Condition} & \textbf{Model} &
\textbf{Mean bias} & \textbf{Mean |bias|} & \textbf{RMSE} \\
\midrule
\multirow{2}{*}{Roleplay} &
  Claude Haiku 4.5 & $+0.650$ & $0.737$ & $1.661$ \\
& Llama 3.1 70B    & $+0.798$ & $0.874$ & $2.024$ \\
\midrule
\multirow{2}{*}{Predict} &
  Claude Haiku 4.5 & $-0.167$ & $0.329$ & $0.329$ \\
& Llama 3.1 70B    & $+0.040$ & $0.189$ & $0.189$ \\
\bottomrule
\end{tabular}
\end{table}

Both models achieve substantially lower RMSE in the predict condition
than in roleplay (Claude: 0.33 vs.\ 1.66; Llama: 0.19 vs.\ 2.02). For
Llama, this reflects genuine accuracy: mean bias of $+0.04$ across all
statistics indicates near-perfect calibration on average. For Claude,
the lower predict RMSE is partly an artifact of cancellation: the mean
bias is $-0.17$, reflecting a systematic reversal of the roleplay
over-production bias, and individual statistics show large under- and
over-production in opposite directions. In roleplay, both models are
net positive (over-producers); in predict, Llama is near zero and
Claude is net negative (under-producer). The models are thus not only
better calibrated in the predict condition---they have qualitatively
different failure modes across conditions.

\subsection{Roleplay Condition: Over-Endogeneity}

Table~\ref{tab:roleplay} reports the statistic-level means and percentage
biases for the roleplay condition. The dominant failure mode is
over-production of history-dependent statistics.

\begin{table}[H]
\centering
\caption{Roleplay condition: mean endogenous statistic values (human empirical
vs.\ LLM-generated), averaged across all domains and networks.
\% bias $= (m_{\text{LLM}} - m_{\text{emp}})/m_{\text{emp}}\times 100$.}
\label{tab:roleplay}
\begin{tabular}{lrrrrrr}
\toprule
\multirow{2}{*}{\textbf{Term}} &
\multirow{2}{*}{\textbf{Empirical}} &
\multicolumn{2}{c}{\textbf{Claude (Roleplay)}} &
\multicolumn{2}{c}{\textbf{Llama (Roleplay)}} \\
\cmidrule(lr){3-4}\cmidrule(lr){5-6}
& & \textbf{Mean} & \textbf{\% bias} & \textbf{Mean} & \textbf{\% bias} \\
\midrule
OutDegSnd      & 2.399 & 6.935 & $+189\%$ & 7.849 & $+227\%$ \\
Repetition     & 0.329 & 0.647 & $+97\%$  & 0.474 & $+44\%$  \\
RecRecencyRank & 0.387 & 0.622 & $+61\%$  & 0.615 & $+59\%$  \\
Imitation      & 0.563 & 0.808 & $+43\%$  & 0.833 & $+48\%$  \\
SndInitiation  & 0.702 & 0.920 & $+31\%$  & 0.929 & $+32\%$  \\
NOTPSnd        & 0.459 & 0.576 & $+25\%$  & 0.767 & $+67\%$  \\
NITPSnd        & 0.410 & 0.482 & $+18\%$  & 0.688 & $+68\%$  \\
WeightedRec    & 0.277 & 0.289 & $+5\%$   & 0.392 & $+42\%$  \\
IOviolation    & 0.732 & 0.544 & $-26\%$  & 0.544 & $-26\%$  \\
\bottomrule
\end{tabular}
\end{table}

\paragraph{OutDegSnd: over-breadth.}
The largest divergence is in sender contact breadth. Human actors contact
a mean of 2.4 normalized distinct partners; Claude reaches 6.9 ($+189\%$)
and Llama 7.8 ($+227\%$). LLMs distribute their interactions across far
more distinct partners than real actors do in any domain.

\paragraph{Repetition and RecRecencyRank: over-concentration on recent history.}
LLMs re-contact their most recent target at dramatically elevated rates
(Repetition: $+97\%$ and $+44\%$) and preferentially respond to being
recently targeted ($+61\%$ for both). Together these biases indicate that
LLMs over-weight the most recent part of the event history: too sticky
(returning too often to their last partner) and too responsive (reacting
too strongly to recent receipt).

\paragraph{IOviolation: under-violation of turn-taking norms.}
IOviolation is the only statistic both LLMs consistently under-produce.
Humans violate sequential participation norms in approximately 73\% of events;
LLMs do so in only 54\% ($-26\%$), with nearly identical deficits for both
models. LLMs are more politely turn-taking than humans and less likely to
break conversational rhythm spontaneously.

\paragraph{WeightedRec: best-calibrated.}
Weighted reciprocity is the most accurately reproduced statistic in the
roleplay condition. Claude is nearly unbiased ($+5\%$); Llama shows
moderate over-weighting ($+42\%$).

\subsection{Predict Condition: Divergent Model Behavior}

Table~\ref{tab:predict} reports the statistic-level results for the predict
condition. The two models diverge markedly.

\begin{table}[H]
\centering
\caption{Predict condition: mean endogenous statistic values (human empirical
vs.\ LLM-generated), averaged across all domains and networks.}
\label{tab:predict}
\begin{tabular}{lrrrrrr}
\toprule
\multirow{2}{*}{\textbf{Term}} &
\multirow{2}{*}{\textbf{Empirical}} &
\multicolumn{2}{c}{\textbf{Claude (Predict)}} &
\multicolumn{2}{c}{\textbf{Llama (Predict)}} \\
\cmidrule(lr){3-4}\cmidrule(lr){5-6}
& & \textbf{Mean} & \textbf{\% bias} & \textbf{Mean} & \textbf{\% bias} \\
\midrule
SndInitiation  & 0.702 & 0.318 & $-55\%$ & 0.799 & $+10\%$ \\
Repetition     & 0.329 & 0.067 & $-80\%$ & 0.410 & $+20\%$ \\
RecRecencyRank & 0.387 & 0.123 & $-72\%$ & 0.432 & $+7\%$  \\
WeightedRec    & 0.277 & 0.129 & $-57\%$ & 0.234 & $-16\%$ \\
Imitation      & 0.563 & 0.309 & $-40\%$ & 0.597 & $+6\%$  \\
OutDegSnd      & 2.399 & 2.158 & $-10\%$ & 2.868 & $+20\%$ \\
NOTPSnd        & 0.459 & 0.456 &  $0\%$  & 0.466 & $+2\%$  \\
NITPSnd        & 0.410 & 0.458 & $+11\%$ & 0.404 & $-1\%$  \\
IOviolation    & 0.732 & 0.904 & $+23\%$ & 0.522 & $-29\%$ \\
\bottomrule
\end{tabular}
\end{table}

\paragraph{Claude in predict: systematic under-production of history-dependent statistics.}
Claude under-produces all four of the robust endogenous baseline statistics
by large margins: SndInitiation ($-55\%$), Repetition ($-80\%$),
RecRecencyRank ($-72\%$), WeightedRec ($-57\%$), and Imitation ($-40\%$).
OutDegSnd, by contrast, is nearly correct ($-10\%$) and triadic closure
terms are close to empirical baselines. IOviolation is over-produced
($+23\%$), the opposite direction from the roleplay condition ($-26\%$).

Claude's predict behavior resembles a relatively uniform or novel-contact
prediction strategy: it forecasts events that ignore established dyadic
relationships, treating the next event as only weakly conditioned on who
has interacted with whom before. The IOviolation reversal is consistent
with this interpretation: when predicting rather than acting,
Claude treats norm violations as the default rather than the exception,
perhaps reflecting a belief that human sequences are more spontaneous
and asymmetric than they actually are.

\paragraph{Llama in predict: near-human calibration.}
Llama reproduces all nine statistics close to the empirical baseline.
The largest deviations are OutDegSnd ($+20\%$), Repetition ($+20\%$),
and IOviolation ($-29\%$). NOTPSnd and NITPSnd are within $2\%$ of
empirical values. This is the best calibration achieved by either model
in either condition and represents a qualitatively different behavior
from Llama's roleplay output, where biases are large and consistent.

\subsection{Condition Comparison: The Roleplay--Predict Reversal}

Table~\ref{tab:compare} places the bias estimates side-by-side to
highlight the roleplay--predict reversal for each model.

\begin{table}[H]
\centering
\caption{Bias comparison across conditions. Cells show
$(m_{\text{LLM}} - m_{\text{emp}})$ averaged across all domains.
Arrows indicate direction of bias: $\uparrow$ = over-production,
$\downarrow$ = under-production.}
\label{tab:compare}
\begin{tabular}{lrrrr}
\toprule
\multirow{2}{*}{\textbf{Term}} &
\multicolumn{2}{c}{\textbf{Claude}} &
\multicolumn{2}{c}{\textbf{Llama}} \\
\cmidrule(lr){2-3}\cmidrule(lr){4-5}
& \textbf{Roleplay} & \textbf{Predict} & \textbf{Roleplay} & \textbf{Predict} \\
\midrule
SndInitiation  & $+0.218\uparrow$ & $-0.409\downarrow$ & $+0.227\uparrow$ & $+0.076\uparrow$ \\
Repetition     & $+0.318\uparrow$ & $-0.337\downarrow$ & $+0.145\uparrow$ & $+0.084\uparrow$ \\
RecRecencyRank & $+0.235\uparrow$ & $-0.309\downarrow$ & $+0.229\uparrow$ & $+0.029\uparrow$ \\
WeightedRec    & $+0.012\uparrow$ & $-0.170\downarrow$ & $+0.116\uparrow$ & $-0.065\downarrow$ \\
Imitation      & $+0.245\uparrow$ & $-0.220\downarrow$ & $+0.269\uparrow$ & $+0.033\uparrow$ \\
OutDegSnd      & $+4.832\uparrow$ & $-0.236\downarrow$ & $+5.893\uparrow$ & $+0.388\uparrow$ \\
NOTPSnd        & $+0.116\uparrow$ & $+0.053\uparrow$   & $+0.307\uparrow$ & $+0.007\uparrow$ \\
NITPSnd        & $+0.072\uparrow$ & $+0.134\uparrow$   & $+0.278\uparrow$ & $-0.008\downarrow$ \\
IOviolation    & $-0.188\downarrow$ & $+0.157\uparrow$ & $-0.214\downarrow$ & $-0.251\downarrow$ \\
\midrule
\textbf{RMSE} & \textbf{1.661} & \textbf{0.329} & \textbf{2.024} & \textbf{0.189} \\
\bottomrule
\end{tabular}
\end{table}

For Claude, every history-dependent statistic reverses sign between
conditions: over-production in roleplay becomes under-production in
predict, and the IOviolation deficit becomes an IOviolation surplus.
The reversal is consistent across all five forward- and reverse-dyad
statistics and across all nine domains. For Llama, the transition from
roleplay to predict does not involve reversals but compressions: all
biases move toward zero.

\subsection{Domain-Level RMSE}

Tables~\ref{tab:rmse_roleplay} and~\ref{tab:rmse_predict} report RMSE
by domain for the roleplay and predict conditions, respectively.

\begin{table}[H]
\centering
\caption{Roleplay condition: RMSE by domain and model.}
\label{tab:rmse_roleplay}
\begin{tabular}{llllrr}
\toprule
\textbf{Domain} & \textbf{Vis.} & \textbf{Val.} &
\textbf{Nets} & \textbf{Claude} & \textbf{Llama} \\
\midrule
War       & High & Conflictual & 1 & 0.134 & 0.836 \\
Email     & Low  & Affiliative & 5 & 0.292 & 0.206 \\
ftf       & Low  & Affiliative & 5 & 0.443 & 0.346 \\
Classroom & High & Affiliative & 1 & 0.600 & 0.444 \\
Gangs     & High & Conflictual & 5 & 1.015 & 1.081 \\
Twitter   & High & Affiliative & 1 & 1.529 & 1.630 \\
PMs       & Low  & Affiliative & 5 & 2.042 & 2.460 \\
Phone     & High & Affiliative & 1 & 2.288 & 2.609 \\
SMS       & Low  & Affiliative & 3 & 7.106 & 8.821 \\
\midrule
\textbf{Overall} & & & 27 & \textbf{1.661} & \textbf{2.024} \\
\bottomrule
\end{tabular}
\end{table}

\begin{table}[H]
\centering
\caption{Predict condition: RMSE by domain and model.}
\label{tab:rmse_predict}
\begin{tabular}{llllrr}
\toprule
\textbf{Domain} & \textbf{Vis.} & \textbf{Val.} &
\textbf{Nets} & \textbf{Claude} & \textbf{Llama} \\
\midrule
ftf       & Low  & Affiliative & 5 & 0.214 & 0.051 \\
Gangs     & High & Conflictual & 5 & 0.198 & 0.258 \\
PMs       & Low  & Affiliative & 5 & 0.230 & 0.160 \\
Twitter   & High & Affiliative & 1 & 0.312 & 0.229 \\
Phone     & High & Affiliative & 1 & 0.364 & 0.138 \\
Email     & Low  & Affiliative & 5 & 0.382 & 0.551 \\
War       & High & Conflictual & 1 & 0.398 & 0.106 \\
Classroom & High & Affiliative & 1 & 0.408 & 0.144 \\
SMS       & Low  & Affiliative & 3 & 0.505 & 0.108 \\
\midrule
\textbf{Overall} & & & 27 & \textbf{0.329} & \textbf{0.189} \\
\bottomrule
\end{tabular}
\end{table}

In the roleplay condition (Table~\ref{tab:rmse_roleplay}), RMSE is
lowest in adversarial or formal settings (War, Email, ftf) and highest
in informal high-volume settings (SMS, Phone, PMs), driven primarily by
OutDegSnd over-production which is most severe in contexts where human
interaction is tightly concentrated on a small number of partners.

In the predict condition (Table~\ref{tab:rmse_predict}), the domain
ordering changes substantially. SMS, which produced the worst roleplay
performance, falls to mid-range for both models (Claude: 0.505; Llama:
0.108) because OutDegSnd is now much more accurately reproduced. Llama
achieves its lowest RMSE in ftf networks (0.051), where human interaction
is tightly structured and the predict task is relatively constrained.
Notably, Email is now one of Llama's \emph{worse} domains (0.551),
reversing its roleplay performance---suggesting that Llama's accurate
prediction of email-like dyadic sequences in roleplay relies partly on
identity assignment rather than pure sequence forecasting.

%------------------------------------------------------------------
\section{Discussion}
\label{sec:discussion}

\subsection{Two Failure Modes, One Contrast}

The comparison between conditions reveals a fundamental asymmetry in how
LLMs represent social interaction under different framing conditions.

In the \emph{roleplay} condition, LLMs over-endogenize: they behave as if
the pull of prior interaction history is stronger than it actually is for
humans. They return to established partners too often, respond too strongly
to recent receipt, engage too many distinct partners, and follow popular
targets too readily. This pattern is consistent with LLMs having internalized
the \emph{structure} of endogenous interaction but having over-fitted to it---
amplifying tendencies that are present in human behavior but exaggerating
their magnitude when adopting a specific agent identity. The pattern
parallels what \citet{cheng2023compost} call \emph{caricature} in
demographic persona simulation: the identity frame activates and exaggerates
the tendencies associated with that identity rather than accurately reproducing
them, producing a heightened version of structural social behavior rather than
a calibrated one. \citet{shanahan2023role}'s concept of \emph{character
capture} provides a complementary account: the roleplay framing can override
baseline model behavior and commit the LLM to acting consistently with the
assigned identity's implied tendencies, even when those tendencies are stronger
than empirically warranted.

In the \emph{predict} condition, Claude reverses course entirely:
it under-produces history-dependent statistics by margins that are as large
as its over-production in roleplay. This is not a calibration improvement; it
is a different error. Claude in predict mode behaves as if social sequences are
\emph{less} structured than they are---as if the next event is only weakly
determined by who has interacted with whom before. Llama, by contrast, shifts
toward the empirical baseline without reversing, suggesting a more stable
underlying representation of sequence statistics.

\subsection{Interpreting Claude's Predict Behavior}

Claude's predict-condition behavior suggests a framing-sensitive
representational shift. When asked to \emph{act} as an agent, Claude
activates a high-inertia social model: return to established partners, follow
network structure, respect turn-taking. When asked to \emph{forecast} what
will happen next in a sequence, Claude appears to activate a different model---
one that discounts dyadic history and treats the next event as relatively
unconditioned on the past. The IOviolation reversal is the clearest marker:
in roleplay, Claude is excessively compliant (under-violating turn-taking
norms); in predict, it forecasts that norm violations are common (over-producing
IOviolation). The empirical truth is between these two representations.

One interpretation is that Claude has two separable implicit models of
social interaction: an \emph{actor model} (used when role-playing an identity)
and a \emph{forecaster model} (used when predicting). The actor model is
high-inertia and norm-compliant; the forecaster model is flatter and
violation-tolerant. Neither is well-calibrated in isolation, but they err
in different directions. This is consistent with \citet{tseng2024two}'s
argument that role-playing and prediction (personalization) are distinct
paradigms with distinct failure modes: conflating them obscures the fact
that what makes an LLM a good actor may be precisely what makes it a poor
forecaster, and vice versa. The coincidental near-accuracy of Claude's
predict RMSE (0.33) is thus misleading: the errors cancel across statistics
rather than reflecting genuine calibration.

\subsection{Interpreting Llama's Predict Accuracy}

Llama's near-perfect calibration in the predict condition (RMSE 0.19,
mean bias $+0.04$) is the most surprising result in this study. Without
any identity assignment or structural scaffolding, Llama reproduces human
interaction statistics nearly exactly on average across nine social domains.

This accuracy may reflect something substantive about Llama's pretraining
data and representational strategy. One possibility is that Llama has
encoded a more accurate \emph{base rate} model of human social interaction
at the sequence level---a statistical summary of who acts and who they target
that is well-calibrated at the population level even if it does not represent
any specific actor's perspective. When asked to forecast the next event,
Llama draws from this base rate model and produces sequences that match the
empirical distribution. When asked to role-play a specific identity, however,
Llama over-exaggerates the structural tendencies implied by that identity,
producing the large positive biases observed in the roleplay condition.

Critically, Llama's predict accuracy should not be taken to mean that Llama
is a better social model overall. In the roleplay condition---where
identity assignment is the realistic framing for agent-based simulation---
Llama's biases are larger than Claude's. The predict accuracy is best
understood as Llama having a well-calibrated \emph{sequence-level prior}
that is decoupled from its \emph{identity-conditioned behavior}.

\subsection{The Visibility and Valence of Bias}

Domain-level RMSE patterns differ markedly between conditions, and the
changes are informative about what each condition is measuring.

In roleplay, the worst performance is in SMS and PMs---low-visibility,
informal, high-volume dyadic settings---primarily because OutDegSnd
over-production is most severe where human actors concentrate most tightly
on a small number of partners. LLMs continue to spread attention broadly
regardless of context.

In predict, the domain ordering flattens dramatically, and SMS is
no longer an outlier. This is because Claude and Llama now choose senders
themselves, and their sender-selection behavior moderates some of the
extreme breadth bias. The removal of REM-governed sender activation
apparently allows more natural variation in who initiates contact, which
pulls OutDegSnd closer to the empirical baseline.

The War domain is the most consistent high performer in roleplay (Claude
RMSE: 0.134) but becomes mid-range in predict for Claude (0.398) and low
for Llama (0.106). This reversal for Claude suggests that the War domain's
strong roleplay performance was substantially driven by the REM-governed
sender activation concentrating events on historically active conflict
dyads, rather than Claude's own calibrated forecasting of conflict dynamics.
Llama's War predict accuracy (0.106) supports a different interpretation:
Llama has genuinely internalized the statistical structure of conflict
sequences well enough to forecast them without scaffolding.

\subsection{Implications for Social Simulation}

The results suggest four practical recommendations for researchers
deploying LLMs in social simulation:

\begin{enumerate}[leftmargin=2em]
  \item \textbf{Condition choice matters as much as model choice.}
        The performance gap between conditions (RMSE 0.19--2.02) is
        larger than the gap between models ($\leq 0.4$ within condition).
        Before selecting a provider, researchers should determine whether
        their application calls for roleplay (identity-based) or predictive
        simulation, as the appropriate calibration target differs
        substantially.

  \item \textbf{Roleplay over-endogeneity requires explicit correction.}
        When using LLMs as identity-assigned agents, OutDegSnd,
        Repetition, and RecRecencyRank will all be substantially
        over-produced. Applications that care about network density,
        partner concentration, or sequence statistics should post-correct
        or re-parameterize the prompt to penalize over-repetition and
        over-breadth.

  \item \textbf{Predict-condition Claude under-produces dyadic structure.}
        For sequence forecasting applications, Claude should not be relied
        upon to reproduce dyadic persistence or reverse-dyad recency
        without additional prompting that explicitly reminds the model to
        weight prior interaction history in its predictions.

  \item \textbf{Llama predict is the strongest off-the-shelf forecaster.}
        For applications requiring aggregate sequence statistics without
        identity assignment, Llama 3.1 70B in the predict condition provides
        the best available calibration across a wide range of social domains.
\end{enumerate}

\subsection{Limitations}

Several limitations bound these conclusions. First, the roleplay
architecture governs sender activation via the empirical REM hazard
surface, which may counteract some LLM-native biases. Purely unconstrained
roleplay (LLM-governed sender activation) might show different, possibly
larger, biases. The roleplay results are thus conservative lower bounds
on the calibration gap under fully endogenous conditions.

Second, both conditions use generic role prompts. Domain-specific prompting
that provides LLMs with detailed sociological context about interaction
norms in each setting might improve calibration, particularly for the
predict condition where no identity frame is provided. Whether the biases
documented here are fundamental to LLM social representations or
prompt-sensitivity artifacts remains an open question.

Third, the evaluation covers nine statistics from one candidate set.
\citet{gibson2025four} find that these nine statistics, plus three others
(DyadRecencyLog, RecRecencyLog, RRecInitiation), capture approximately 97\%
of the AIC improvement achievable from a twelve-term endogenous candidate
set---but LLM biases on terms not represented here cannot be ruled out.

Fourth, Llama's predict accuracy---while genuine---may not generalize to
settings or domains not represented in this corpus. The nine domains here
are diverse, but they all involve human social interaction in Anglophone
contexts, and they are all represented in the kind of text data on which
large language models are typically trained.

%------------------------------------------------------------------
\section{Conclusion}

We have shown that LLMs behave differently as social actors than as social
forecasters, and that the two behaviors err in systematically different
directions. In the roleplay condition, both Claude Haiku 4.5 and Llama 3.1
70B over-produce history-dependent interaction statistics---too much
repetition, recency-seeking, contact breadth, and imitation---while
under-producing the spontaneous norm violations that characterize real
human interaction. In the predict condition, Claude reverses these biases,
sharply under-producing the same history-dependent statistics, while Llama
converges on the empirical baseline with near-human accuracy.

These findings resist a simple ``more scaffolding is better'' or ``predict
is harder'' interpretation. Llama's predict accuracy is not an artifact of
the task being easy; it reflects a well-calibrated implicit model of
sequence-level social dynamics that is decoupled from the over-exaggeration
of identity-conditioned behavior in roleplay. Claude, by contrast, has two
poorly calibrated but internally coherent models---an actor model that
over-structures behavior and a forecaster model that under-structures it---
neither of which is close to correct.

For the growing community of researchers deploying LLMs as social simulation
agents, these results argue for treating the interaction between model choice
and condition framing as a primary design variable, not a secondary concern.
The same LLM that produces acceptable results in one framing may be
severely miscalibrated in another. Before treating LLM-generated interaction
sequences as valid proxies for human behavior, empirical calibration across
the endogenous statistics that structure real social interaction is warranted.

%------------------------------------------------------------------
\section*{Acknowledgements}

\placeholder{Insert acknowledgements here.}

\section*{Conflicts of Interest}

The author declares no conflicts of interest.

\section*{Data Availability Statement}

The empirical sequences used in this study are drawn from multiple
previously published corpora. \placeholder{Specify dataset availability.}
Analysis code and generated sequences are available at
\url{https://github.com/\ldots}.

%------------------------------------------------------------------
\bibliography{cites}
\bibliographystyle{plainnat}

\end{document}
